\documentclass[11pt, letterpaper]{article}

\usepackage[margin=0.7in]{geometry}
\usepackage{hyperref}
\usepackage{booktabs}
\usepackage{longtable}
\usepackage{parskip}
\usepackage{titlesec}
\usepackage{enumitem}
\usepackage{graphicx}
\usepackage{float}
\usepackage{subcaption}
\usepackage{caption}
\usepackage{amsmath}
\usepackage{amssymb}
\usepackage{array}

\titleformat{\section}{\large\bfseries}{}{0em}{}[\titlerule]
\titleformat{\subsection}{\normalsize\bfseries}{}{0em}{}

\hypersetup{
    colorlinks=true,
    linkcolor=black,
    urlcolor=blue,
    citecolor=black
}

\graphicspath{{analysis/}}

\title{
    \textbf{Progress Report} \\[0.4em]
    \large Architecture and Training Design Under a Hard Parameter Budget \\[0.4em]
    \normalsize DS 504: Big Data Analytics \textendash{} Spring 2026
}

\author{Shyam Patadia \\ \\ Worcester Polytechnic Institute}
\date{April 21, 2026}

\begin{document}

\maketitle
\thispagestyle{empty}

% -------------------------------------------------------
\section{Summary of Progress}
% -------------------------------------------------------

This report documents progress on the controlled ablation study
of architectural and training design choices for decoder-only
language models under the 16~MB artifact constraint of the
OpenAI Model Craft: Parameter Golf Challenge~\cite{paramgolf2026}.
Since the proposal, all three primary objectives have reached
convergent 20{,}000-step runs at seed 1337 on the WPI research
cluster: O1 (quantization precision), O2 (positional encoding),
and O3 (optimizer). An additional fourth exploratory axis --
parameter-efficient weight factorizations (low-rank, block-diagonal,
fixed random projection) -- was evaluated and falsified.
Combining the best per-axis configurations into a single
stacked run exposed a capacity-versus-quantization-noise
tradeoff that is the central finding of this phase and directly
motivates the remaining work.

Key result: the best in-axis condition (partial RoPE, 9L/MLP=2,
INT8) reaches 1.2147~bpb, a 0.010~bpb improvement over the
1.2244~bpb reference. A stacked configuration
(partial RoPE + INT6 + 11~layers + MLP multiplier~3) reaches
1.1614~bpb during training but degrades to 1.2816~bpb post-quantization
-- a 0.14~bpb gap that identifies INT6 uniform quantization as
the dominant source of end-to-end loss in the capacity-expanded
regime. Remaining work decomposes this gap along two orthogonal axes:
a TurboQuant random-rotation quantizer~\cite{turboquant2024}
applied to the same expanded architecture, and a pure
architectural expansion (11L) under the known-safe INT8
quantizer. Both jobs are in flight.

% -------------------------------------------------------
\section{Experimental Setup}
% -------------------------------------------------------

\subsection{Dataset and Tokenization}

All models are trained from scratch on FineWeb~\cite{fineweb2024},
a filtered, deduplicated Common Crawl derivative released by
HuggingFace. Training uses a 10~billion-token shard
(\texttt{fineweb10B}) and the 1{,}024-token SentencePiece BPE
tokenizer (\texttt{sp1024}) adopted by the Parameter Golf
challenge; the small vocabulary keeps the embedding table to
$1024 \times d_{\text{model}}$ bytes under FP32, which matters
substantially under a 16~MB artifact budget. Input/output
embeddings are tied throughout.

Training reads contiguous token streams of length
$\ell_{\text{ctx}}=1024$ directly from pre-shuffled, pre-tokenized
\texttt{uint16} binary shards; no on-the-fly tokenization is
performed during training. Evaluation computes bits-per-byte
(bpb) on a held-out FineWeb validation split using the official
challenge evaluation protocol, so that our numbers are directly
comparable to leaderboard submissions~\cite{paramgolf2026}.
Total training exposure per condition is $\approx 10.5$~billion
tokens (20{,}000 steps $\times$ 8 micro-batches $\times$
2 GPUs $\times$ 1024 tokens, adjusted for gradient accumulation).

\subsection{Model Architecture}

The baseline architecture is a decoder-only transformer with the
hyperparameters in Table~\ref{tab:arch}. Each block applies
pre-norm $\rightarrow$ causal multi-head self-attention
$\rightarrow$ pre-norm $\rightarrow$ gated MLP (SwiGLU), with a
residual stream of width $d_{\text{model}}$. Ablations mutate
exactly one axis per run relative to the baseline; arch-expansion
runs (O5) mutate several axes jointly within the same artifact
budget.

\begin{table}[H]
\centering
\small
\begin{tabular}{@{}lll@{}}
\toprule
\textbf{Component} & \textbf{Symbol} & \textbf{Baseline value} \\
\midrule
Layers                    & $L$                  & $9$ \\
Hidden dimension          & $d_{\text{model}}$   & $512$ \\
Attention heads           & $h$                  & $8$ \\
Head dimension            & $d_h$                & $64$ \\
MLP expansion             & $m$                  & $2$ (SwiGLU) \\
Context length            & $\ell_{\text{ctx}}$  & $1024$ \\
Vocabulary                & $|V|$                & $1024$ (SentencePiece BPE) \\
Tied embeddings           & --                   & input $=$ output \\
Training steps            & $T$                  & $20{,}000$ \\
Micro-batch / GPU / accum & --                   & $8 / 2 / \text{as-needed}$ \\
Optimizer (baseline)      & --                   & Muon~\cite{muon2024} on matrix params, Adam on scalars \\
Quantization (baseline)   & --                   & INT8 PTQ with per-row scales, \texttt{zlib} roundtrip \\
\bottomrule
\end{tabular}
\caption{Baseline architecture and training hyperparameters. All
experiments mutate exactly one axis from this baseline unless
otherwise noted.}
\label{tab:arch}
\end{table}

\subsection{Quantization and Artifact Pipeline}

The 16~MB budget applies to the \emph{post-quantization, zlib-compressed}
artifact plus training code. To mirror the leaderboard
evaluation, every reported run performs the following
artifact-production and round-trip pipeline at the end of
training:

\begin{enumerate}[leftmargin=*, topsep=0pt, itemsep=1pt]
    \item Train in bfloat16 (or with QAT straight-through
          estimators for INT6).
    \item Quantize weights to the target precision with per-row
          symmetric scales; keep 1D tensors (biases, norms) in
          bfloat16 to avoid disproportionate error.
    \item Serialize the quantized state-dict and \texttt{zlib.compress}
          it; this is the artifact measured against 16~MB.
    \item \emph{Round-trip}: decompress, dequantize, rehydrate
          the model, re-run validation. This is the leaderboard-equivalent
          \texttt{final/val\_bpb} reported throughout.
\end{enumerate}

Two distinct bpb numbers therefore exist for every run:
\texttt{val/bpb} (training-time, bfloat16 in-memory parameters)
and \texttt{final/val\_bpb} (post round-trip, what the leaderboard
actually measures). The gap between them is the quantization
penalty, and Section~\ref{sec:results} shows it is the central
empirical object of this phase.

% -------------------------------------------------------
\section{Results by Objective}
\label{sec:results}
% -------------------------------------------------------

Table~\ref{tab:results} aggregates the full grid of 20{,}000-step
runs. All numbers are at seed 1337; multi-seed replication
(seeds 42 and 123) is the final-phase deliverable. The
\textsc{Final} column is leaderboard-equivalent bpb; the
\textsc{Train} column is the in-memory bfloat16 validation
bpb at the end of training. Their difference is the
quantization gap $\Delta_q$.

\begin{table}[H]
\centering
\small
\begin{tabular}{@{}l l c c c c@{}}
\toprule
\textbf{Obj.} & \textbf{Condition} & \textbf{Train bpb}
              & \textbf{Final bpb} & $\Delta_q$
              & \textbf{vs.\ 1.2244} \\
\midrule
O1 & INT8 PTQ (baseline)              & 1.2082 & 1.2161 & $+0.008$ & $-0.008$ \\
O1 & INT6 QAT + PTQ                   & 1.2083 & 1.3489 & $+0.141$ & $+0.125$ \\
\midrule
O2 & Learned absolute PE              & 1.2244 & 1.2349 & $+0.011$ & $+0.011$ \\
O2 & Full RoPE                        & 1.2085 & 1.2174 & $+0.009$ & $-0.007$ \\
O2 & No positional encoding           & 1.2589 & 1.2690 & $+0.010$ & $+0.045$ \\
\textbf{O2} & \textbf{Partial RoPE ($d_{\text{rope}}{=}32$)}
                                      & \textbf{1.2071} & \textbf{1.2147}
                                      & $+0.008$ & $\mathbf{-0.010}$ \\
\midrule
O3 & AdamW (5k pilot)                 & 1.4065 & 1.4409 & $+0.034$ & $+0.217$ \\
O3 & Adafactor (5k pilot)             & 1.4879 & 1.4895 & $+0.002$ & $+0.265$ \\
O3 & Muon (baseline, 20k)             & 1.2089 & 1.2168 & $+0.008$ & $-0.008$ \\
\midrule
O4 & Low-rank $r{=}128$               & 1.2915 & 1.3014 & $+0.010$ & $+0.077$ \\
O4 & Block-diag $g{=}4$               & 1.3700 & 1.3830 & $+0.013$ & $+0.159$ \\
O4 & Fixed random proj.\ $k{=}128$    & 1.3838 & 1.3958 & $+0.012$ & $+0.171$ \\
\midrule
O5 & Partial RoPE + INT6 + 11L + MLP=3 & \textbf{1.1614} & 1.2816 & $\mathbf{+0.120}$ & $+0.057$ \\
\bottomrule
\end{tabular}
\caption{Full 20{,}000-step results at seed 1337, reported as
bits-per-byte on the FineWeb held-out split. \textsc{Train} is
the in-memory bfloat16 loss; \textsc{Final} is the post-quantize,
zlib round-trip loss (leaderboard-equivalent). $\Delta_q$ is the
quantization gap. Reference 1.2244 is the challenge's published
INT8 score for this architecture.}
\label{tab:results}
\end{table}

\subsection{O1 -- Quantization Precision}

Figure~\ref{fig:o1} plots validation bpb against step for the
INT8 and INT6 conditions. In-memory training loss is nearly
identical between the two (both $\approx 1.208$ at step 20{,}000);
the critical difference appears only after the round-trip.
INT8 with per-row symmetric scales incurs a negligible
$\Delta_q = 0.008$~bpb. INT6 under the same scheme incurs
$\Delta_q = 0.141$~bpb -- more than sixteen times as large
-- even with QAT applied throughout. The INT6 penalty is not a
training-dynamics problem: the straight-through estimator reaches
the same loss as INT8, but uniform scalar quantization at 6~bits
cannot represent the post-training weight distribution without
appreciable error. This result contradicts the optimistic
prediction from existing precision scaling
laws~\cite{scalingprecision} in the sub-21M parameter regime
and reframes INT6 as a capacity multiplier that is only useful
if a better quantizer is supplied.

\begin{figure}[H]
\centering
\includegraphics[width=0.80\textwidth]{curves_o1_quant.png}
\caption{O1 -- post-warmup validation bpb over training steps
for INT8 (baseline) and INT6 (QAT). In-memory loss converges
to the same point; the divergence shown in
Table~\ref{tab:results} emerges only under the post-training
round-trip.}
\label{fig:o1}
\end{figure}

\subsection{O2 -- Positional Encoding}

Four positional-encoding schemes were compared under
otherwise-identical architecture and quantization
(INT8 PTQ, 9L, MLP=2). The ordering at 20{,}000 steps is
\emph{partial RoPE $<$ full RoPE $<$ learned absolute $\ll$ none},
with partial RoPE at 1.2147~bpb becoming the best in-axis
condition of this study. The gap between partial RoPE and no-PE
(1.2690) is 0.054~bpb, dwarfing the gap between partial RoPE
and full RoPE (0.003~bpb). Absolute-position information matters
at this scale; the \emph{form} of the positional prior is a
second-order variable. Partial RoPE's advantage specifically
comes from applying rotary encoding to only the first 32 of
each head's 64 dimensions, leaving the remainder as pure
content dimensions and freeing parameters that learned
embeddings would otherwise require~\cite{rope}.

\begin{figure}[H]
\centering
\includegraphics[width=0.80\textwidth]{curves_o2_posenc.png}
\caption{O2 -- validation bpb over training for the four
positional-encoding variants. \emph{No PE} trails clearly;
the three schemes with position information converge to
within 0.02~bpb of each other, with partial RoPE lowest.}
\label{fig:o2}
\end{figure}

\subsection{O3 -- Optimizer}

The two comparison optimizers are included for distinct,
budget-relevant reasons. \textbf{AdamW} is the universal default
in language-model training; without an AdamW counterfactual,
the choice of Muon would be an unjustified pick rather than a
finding, and the central claim that ``Muon is load-bearing at
this scale'' would be unfalsifiable. \textbf{Adafactor} is
specifically motivated by the 16~MB constraint: its factorized
second moment reduces optimizer-state memory from $O(n)$ to
$O(\sqrt{n})$, which directly tests whether sacrificing
second-moment fidelity buys back enough training-time GPU
memory to fit a deeper or wider model under the same wallclock.
Both are therefore lever-tests, not throat-clearing.

The data falsifies both alternatives. Muon, which applies
Nesterov-momentum updates along the Newton--Schulz-orthogonalized
gradient of each 2D weight matrix~\cite{muon2024}, reaches
1.2168~bpb and is indistinguishable from the challenge-reference
INT8 baseline. AdamW at tuned learning rates plateaus $>$0.2~bpb
higher; Adafactor's factorized second moment loses even more
(see Table~\ref{tab:results}). At under 5M trainable parameters,
the optimizer-state savings Adafactor offers do not compensate
for its loss-quality penalty, and AdamW's per-parameter second
moment provides no advantage over Muon's orthogonalization at
this scale. Muon is therefore retained as the matrix-parameter
optimizer for all subsequent experiments, with Adam on scalars.

\begin{figure}[H]
\centering
\includegraphics[width=0.80\textwidth]{curves_o3_optim.png}
\caption{O3 -- validation bpb over training for the optimizer
comparison. Muon converges to a qualitatively different loss
band than either AdamW or Adafactor at this scale.}
\label{fig:o3}
\end{figure}

\subsection{O4 -- Parameter-Efficient Factorizations (exploratory)}

Three structured weight parametrizations were evaluated to
test whether they could liberate capacity under the 16~MB
budget: low-rank ($W = UV^\top$ with $r = 128$), block-diagonal
($g = 4$), and fixed random projection ($W = PR$ with frozen
random $R$, rank 128). All three landed between 1.30 and 1.40~bpb,
cleanly worse than the dense INT8 baseline at 1.2161. The random-projection
variant was further retested at 12 layers to rule out capacity
starvation; it improved to 1.3601 but still failed to clear
baseline. O4 is therefore reported as a null result: under this
budget and token horizon, structured weight factorizations
cannot compete with dense matrices. The implication is that
capacity expansion, if desired, should come from depth/width
rather than from a parametrization trick.

\begin{figure}[H]
\centering
\includegraphics[width=0.80\textwidth]{curves_o4_weightparam.png}
\caption{O4 -- weight-parametrization variants all converge
above the dense INT8 baseline, regardless of which structural
prior is imposed. This axis is falsified as a capacity lever
at our scale.}
\label{fig:o4}
\end{figure}

\subsection{O5 -- Stacked Configuration (the central finding)}

Combining O1--O3's winning choices (partial RoPE, INT6 for
capacity, Muon) with a widened architecture (11 layers,
MLP multiplier 3) produces the most informative run of this
phase. In training-time bfloat16, the stacked configuration
reaches \textbf{1.1614~bpb} -- 0.047~bpb below the best
single-axis result and 0.063~bpb below the reference 1.2244.
After the round-trip, however, it degrades to 1.2816~bpb,
forfeiting the entire training-time gain and landing \emph{worse}
than the INT8 baseline.

The decomposition is unambiguous. The 11L/MLP=3 arch expansion
genuinely improves in-memory capacity; the INT6 precision
chosen to fund that expansion imposes a 0.12~bpb quantization
penalty (consistent with the 0.14~bpb penalty measured
independently in O1). The two effects nearly cancel. This is
the key empirical finding of the phase: \emph{at this scale,
post-quantization loss -- not training-time loss -- is the
binding constraint, and the dominant knob governing post-quant
loss is the quantizer itself, not the architecture}.

\begin{figure}[H]
\centering
\includegraphics[width=0.80\textwidth]{curves_o5_stacked.png}
\caption{O5 -- the stacked configuration (partial RoPE + INT6
+ 11L + MLP=3) vs.\ the INT8 baseline. Training-time loss
(dashed) shows a clear capacity win; final round-trip loss
(solid) does not. The gap between the two lines is the
quantization penalty.}
\label{fig:o5}
\end{figure}

% -------------------------------------------------------
\section{Key Findings}
% -------------------------------------------------------

\begin{enumerate}[leftmargin=*]
    \item \textbf{INT6 uniform scalar quantization is fragile at
          this scale.} A $\Delta_q$ of 0.14~bpb dwarfs every
          other axis's effect. QAT does not close this gap;
          training-time and in-memory losses are already matched.
          The problem is the quantizer, not the training recipe.
    \item \textbf{Partial RoPE ($d_{\text{rope}}=32$) is the
          strongest in-axis modification.} 1.2147~bpb post-round-trip
          beats all other single-axis conditions and equals the
          stacked configuration's \emph{final} loss without any
          of the quantization risk.
    \item \textbf{Muon is load-bearing, not marginal.} At under
          5M parameters, AdamW and Adafactor lose more than 0.2~bpb.
          This axis should not be a free parameter in subsequent
          experiments.
    \item \textbf{Structured weight factorizations do not help.}
          O4 is a falsified hypothesis and removes a large
          region of the design space from consideration.
    \item \textbf{Capacity expansion is real but quantization-gated.}
          The stacked configuration demonstrates a
          $\approx$~0.06~bpb training-time gain from 11L/MLP=3,
          which is recoverable if and only if we can deliver
          a 3--4-bit quantizer with sub-0.02~bpb round-trip
          penalty. That is the remit of the final phase.
\end{enumerate}

% -------------------------------------------------------
\section{In-Progress and Planned Work}
% -------------------------------------------------------

\subsection{TurboQuant at 4~bits (in flight)}

The INT6 fragility result directly motivates TurboQuant~\cite{turboquant2024},
which attacks exactly the failure mode we observe. TurboQuant
applies a random rotation $R$ to each matrix row before scalar
quantization; because a random rotation of an arbitrary
$\ell_2$-bounded vector is concentrated on a Beta distribution
with known moments, a Lloyd--Max scalar quantizer designed for
that distribution achieves provably near-optimal MSE at very
low bit widths. TurboQuant reports quality neutrality at
3.5~bits on LLaMA KV-cache quantization~\cite{turboquant2024};
we apply the same machinery to model \emph{weights} in the
post-training quantizer at 4~bits, via the pure-PyTorch
\texttt{pyturboquant} reference implementation.

Two jobs are currently queued, addressing the two distinct
ways the post-quantization gap can be closed:
\begin{itemize}[leftmargin=*, topsep=0pt, itemsep=1pt]
    \item \textbf{TurboQuant-INT4 at 11L + MLP=3}: a direct
          rerun of the stacked configuration with INT6 replaced
          by TurboQuant-INT4. This isolates the quantizer as
          the variable: same architecture, same token budget,
          same in-memory capacity target, only the
          quantization scheme changes. If TurboQuant reduces
          $\Delta_q$ to the $<0.02$~bpb range observed for INT8,
          the training-time 1.1614 carries through to the
          final score.
    \item \textbf{INT8 at 11L + MLP=2}: keeps INT8 PTQ fixed and
          expands capacity by one architectural notch above the
          baseline (9L $\rightarrow$ 11L). This isolates
          architecture as the variable under a quantizer with
          known, small $\Delta_q$, and quantifies how much of
          the stacked configuration's training-time gain is
          attributable to depth alone versus the
          depth\,$+$\,MLP-width combination.
\end{itemize}

Together the two jobs decompose the O5 result along orthogonal
axes: one varies the quantizer with architecture held at the
expanded setting, the other varies the architecture with the
quantizer held at the safe setting. The combined evidence
identifies which lever -- precision or capacity -- is the
binding constraint at this budget.

\subsection{Per-Condition Learning-Rate and Warmup Sweep}

The proposal specified that learning rate and warmup be tuned
independently per condition. To stay within cluster budget while
covering the full design grid, this phase used a single
hyperparameter setting (the baseline's tuned values) across all
ablations, treating it as a fair-comparison shared baseline
rather than per-condition optima. This is a known source of
under-reporting: at this scale a proper LR/warmup sweep is
typically worth 0.01--0.03~bpb per condition, and conditions
that change effective gradient scale most (e.g., the partial
RoPE attention modification, MLP-multiplier changes, and the
TurboQuant-INT4 weight distribution) are the most likely to
benefit from re-tuning.

The remaining phase will execute a per-condition sweep over
$\{0.5, 1.0, 2.0\}\times$ the baseline learning rate and over
warmup fractions $\{0.05, 0.10, 0.15\}$ for the configurations
that survive into the multi-seed phase. Reported final numbers
for those configurations will use the per-condition optimum
rather than the shared baseline setting; the difference will
be itemized in the final report so the contribution of tuning
versus the contribution of each ablation is visible.

\subsection{Multi-Seed Replication}

The proposal's three-seeds-per-condition protocol will be
executed for the configuration that yields the lowest final
bpb among the two in-flight jobs and for the core axis
ablations (partial RoPE, full RoPE, learned, none; INT8 vs.\
INT6; Muon), each at its per-condition LR/warmup optimum from
the sweep above. Seeds 42 and 123 are already queued in
\texttt{next\_steps.sh} and will be launched as cluster
bandwidth allows. Final reported numbers will be
mean $\pm$ standard deviation across three seeds.

\subsection{Tentative Demonstration}

Scope permitting, the final 16~MB artifact will be packaged
as an in-browser interpretability tool: a static web page that
accepts a user-typed prefix and renders, in real time, the
top-$k$ next-token distribution with probability bars, the
per-head attention heatmap over input tokens, and the
hidden-state magnitudes across layers. Inference runs entirely
client-side via ONNX Runtime Web; no server is involved.
The choice of an interpretability surface rather than a generic
autocomplete UI is deliberate -- the model's small size makes
every internal state cheap to render live, which is infeasible
at the scale of conventional language models, and turns the
16~MB constraint into a pedagogical feature rather than just a
scoring axis. This is a stretch deliverable; the report's
quantitative results stand independently of whether it ships.

\subsection{Revised Schedule}

\begin{longtable}{@{}p{2.8cm} p{9cm} p{2.8cm}@{}}
\toprule
\textbf{Period} & \textbf{Tasks} & \textbf{Deliverable} \\
\midrule
\endfirsthead
Apr 21 -- Apr 27 &
Complete TurboQuant-INT4 and INT8 11L-expansion 20k runs;
attribute the O5 training-time gain to quantizer vs.\
architecture; per-condition LR/warmup sweep on surviving
configurations &
Quantizer/architecture attribution; tuned hyperparameters \\[4pt]
Apr 28 -- May 3 &
Multi-seed replication at per-condition LR optima;
mean $\pm$ std tables; isolate the tuning contribution
from the ablation contribution; write final report;
prepare presentation &
Camera-ready report \\[4pt]
\textbf{May 5} &
\textbf{Final project due} &
\textbf{Report, code, slides} \\
\bottomrule
\end{longtable}

% -------------------------------------------------------
\section{Technical Stack}
% -------------------------------------------------------

Python 3.12, PyTorch 2.8 (CUDA 12.6), HuggingFace
\texttt{datasets} and \texttt{tokenizers}, \texttt{torchao}
for the INT6 QAT path, \texttt{pyturboquant}~\cite{turboquant2024}
for random-rotation post-training quantization, Muon as the
matrix-parameter optimizer, Weights \& Biases for experiment
tracking, SLURM for cluster job scheduling on the WPI academic
partition (2$\times$GPU per job).

% -------------------------------------------------------
\section*{References}
\begingroup
\renewcommand{\section}[2]{}
\begin{thebibliography}{9}

\bibitem{paramgolf2026}
OpenAI.
\newblock {OpenAI Model Craft: Parameter Golf Challenge}, 2026.
\newblock \url{https://github.com/openai/parameter-golf}

\bibitem{fineweb2024}
Penedo, G. et al.
\newblock {The FineWeb Datasets: Decanting the Web for the Finest
Text Data at Scale}.
\newblock \textit{arXiv preprint arXiv:2406.17557}, 2024.

\bibitem{scalingprecision}
Kumar, K. et al.
\newblock {Scaling Laws for Precision}.
\newblock \textit{ICLR}, 2025.
\newblock \textit{arXiv preprint arXiv:2403.08540}.

\bibitem{rope}
Su, J. et al.
\newblock {RoFormer: Enhanced Transformer with Rotary Position
Embedding}.
\newblock \textit{arXiv preprint arXiv:2104.09864}, 2021.

\bibitem{muon2024}
Jordan, K. et al.
\newblock {Muon: An Optimizer for the Hidden Layers of Neural
Networks}.
\newblock Technical report, 2024.

\bibitem{turboquant2024}
Zandieh, A. et al.
\newblock {TurboQuant: Online Vector Quantization with Near-Optimal
Distortion Rate}.
\newblock \textit{arXiv preprint}, 2024.

\bibitem{kaplan2020}
Kaplan, J. et al.
\newblock {Scaling Laws for Neural Language Models}.
\newblock \textit{arXiv preprint arXiv:2001.08361}, 2020.

\bibitem{hoffmann2022}
Hoffmann, J. et al.
\newblock {Training Compute-Optimal Large Language Models}.
\newblock \textit{arXiv preprint arXiv:2203.15556}, 2022.

\end{thebibliography}
\endgroup

\end{document}
